\chapter*{Dédicaces}
\addcontentsline{toc}{chapter}{Dédicaces}
\frontmark{Dédicaces}

\vspace*{0.3cm}
\begin{center}
\itshape
À nos parents,\\
pour leur amour, leur patience\\
et leur soutien à chaque étape de nos études.\\[1.4em]
À nos familles et à nos amis,\\
qui ont cru en nous, même pendant les nuits de calcul.\\[1.4em]
À nos enseignants,\\
qui nous ont appris à poser les bonnes questions\\
avant de chercher les réponses.
\end{center}

\vspace{1.2cm}
{\normalfont\bfseries\fontsize{24}{28}\selectfont Remerciements\par}
\addcontentsline{toc}{chapter}{Remerciements}
\vspace{0.6cm}

Nous remercions très sincèrement \textbf{M. Fayçal Ben Hmida}, enseignant du module « Sûreté de fonctionnement et maintenance », pour la clarté de son cours, pour ses conseils pendant les séances de suivi et pour les points précis qu'il nous a demandé de traiter : ils ont guidé toute la structure de ce rapport.

Nous remercions aussi l'ensemble des enseignants de l'École Nationale Supérieure d'Ingénieurs de Tunis (ENSIT), qui nous ont donné les outils de statistique, de fiabilité et d'apprentissage automatique utilisés ici.

Merci enfin à l'institut FEMTO-ST et aux organisateurs du challenge IEEE PHM 2012, qui ont rendu publiques les données du banc PRONOSTIA, sans lesquelles ce travail n'aurait pas été possible.

{\setstretch{1.0}\small\tableofcontents}

\clearpage\frontmark{Liste des figures et des tableaux}
{\setstretch{1.0}\small
\listoffigures
\begingroup\let\clearpage\relax\vspace{-1.3cm}\listoftables\endgroup
}

\chapter*{Liste des abréviations}
\addcontentsline{toc}{chapter}{Liste des abréviations}
\frontmark{Liste des abréviations}

{\setstretch{1.05}\small\setlength{\parskip}{2.5pt}
\textbf{AFT} : modèle à temps de vie accéléré (\emph{accelerated failure time}).

\textbf{AIC} : critère d'information d'Akaike (plus petit = meilleur compromis entre précision et simplicité).

\textbf{B10, B50} : âge auquel 10 \% (resp. 50 \%) des roulements ont défailli.

\textbf{CBM} : maintenance conditionnelle (\emph{condition-based maintenance}).

\textbf{CNN 1D} : réseau de neurones à convolution sur une dimension (le temps).

\textbf{FPT} : instant du premier défaut détecté (\emph{first predicting time}).

\textbf{GPR} : régression par processus gaussien.

\textbf{HI} : indicateur de santé (\emph{health indicator}).

\textbf{IC} : intervalle de confiance.

\textbf{KPI} : indicateur clé de performance (\emph{key performance indicator}).

\textbf{L10} : durée de vie nominale des catalogues (atteinte par 90 \% des roulements).

\textbf{LSTM} : réseau de neurones récurrent à mémoire longue (\emph{long short-term memory}).

\textbf{MAE} : erreur absolue moyenne.

\textbf{MRL} : durée de vie résiduelle moyenne (\emph{mean residual life}).

\textbf{MTTF} : durée moyenne de fonctionnement avant défaillance (\emph{mean time to failure}).

\textbf{MV} : maximum de vraisemblance.

\textbf{PHM} : surveillance de la santé et pronostic (\emph{prognostics and health management}).

\textbf{RF} : forêt aléatoire (\emph{random forest}).

\textbf{RMSE} : racine de l'erreur quadratique moyenne.

\textbf{RUL} : durée de vie résiduelle (\emph{remaining useful life}).

\textbf{SVR} : régression à vecteurs de support (\emph{support vector regression}).

\textbf{XAI} : intelligence artificielle explicable (\emph{explainable AI}).
}
